\documentclass[11pt]{article}

\usepackage[margin=1in]{geometry}
\usepackage[T1]{fontenc}
\usepackage[utf8]{inputenc}
\usepackage{amsmath}
\usepackage{amssymb}
\usepackage{array}
\usepackage{booktabs}
\usepackage{enumitem}
\usepackage{hyperref}
\usepackage{longtable}
\usepackage{xurl}

\hypersetup{colorlinks=true, linkcolor=blue, urlcolor=blue, citecolor=blue}
\urlstyle{same}
\setlist{nosep,leftmargin=*}
\setlength{\parindent}{0pt}
\setlength{\parskip}{0.65em}

\title{Current SpaceAGORA\_RL Implementation}
\author{SpaceAGORA Contributors}
\date{\today}

\begin{document}
\maketitle

\begin{abstract}
This document summarizes the current Julia implementation of
\path{SpaceAGORA_RL.jl}.  The implementation provides a paper-compatible
aerobraking Markov decision process (MDP), a Julia-native Double Deep
Q-Network (DDQN) learner, baseline policies, evaluation utilities, and
threaded rollout orchestration.  At present, the default and only supported
simulation backend mode is the deterministic \path{:paper_surrogate} backend;
the full SpaceAGORA propagation adapter is represented by an interface hook
but is not yet implemented.
\end{abstract}

\section{Scope and Source Layout}

The RL package lives under \path{SpaceAGORA_RL.jl}.  Its top-level module
includes five implementation areas:

\begin{center}
\begin{tabular}{>{\raggedright\arraybackslash}p{0.32\linewidth}
                >{\raggedright\arraybackslash}p{0.58\linewidth}}
\toprule
Area & Responsibility \\
\midrule
\path{src/core} &
Abstract interfaces for MDPs, simulation backends, policies, transitions, and
episode summaries. \\
\path{src/tasks/aerobraking} &
The Mars Odyssey aerobraking task: observations, normalization, actions,
reward, termination, randomization, backend dynamics, baselines, and
evaluation. \\
\path{src/algorithms/ddqn} &
Replay buffer, Q-network, epsilon schedule, target-network update logic,
DDQN learner, checkpoints, and greedy checkpoint policies. \\
\path{src/runtime} &
Configuration loading, run manifests, and training-session construction. \\
\path{src/parallel} &
Threaded rollout protocol and central learner ingestion.  A small
\path{Distributed} helper exists, but the active training path is
thread-based. \\
\bottomrule
\end{tabular}
\end{center}

Reference Python material remains under \path{Reference Code/}, but it is not
included on the Julia runtime path.  Generated run artifacts are written under
\path{outputs/}.

\section{Core Interfaces}

The package uses small function-based contracts rather than a large framework.
The core abstractions are:

\begin{itemize}
\item \path{AbstractMDP}: task-level marker type.
\item \path{AbstractRLBackend}: reset and step interface for environment
execution.
\item \path{AbstractPolicy}: policy interface returning a 1-based discrete
action index.
\item \path{Transition}: learner-facing tuple containing normalized
observation, action index, reward, next normalized observation, termination
flags, and an info index.
\item \path{EpisodeSummary}: evaluation-facing aggregate containing pass count,
reward, success and failure flags, target error, mission duration, total
delta-V, maneuver count, solver failures, and per-pass traces.
\end{itemize}

The important generic functions are:

\begin{itemize}
\item \path{reset_scenario(config, rng)} returns the initial decision state.
\item \path{step_scenario(config, state, action, rng)} returns an
\path{AerobrakingStepResult}.
\item \path{observe_state(config, state)} returns a \path{PaperObservation}.
\item \path{normalize_observation(obs, bounds)} returns a 9-vector of
learner-facing features.
\item \path{policy_action_index(policy, config, state, observation, rng)}
returns an action index in $\{1,\ldots,13\}$.
\end{itemize}

\section{Aerobraking Task}

\subsection{Scenario Configuration}

The implemented task is a Mars Odyssey aerobraking MDP.  The scenario
configuration stores the phase constants, Mars constants, thermal model
constants, backend mode, reward and termination settings, randomization
settings, normalization bounds, and a training-mode flag.

The helper \path{mars_odyssey_phase_constants} currently supports the
following phases:

\begin{center}
\begin{tabular}{lrrrr}
\toprule
Phase & Target $r_a$ (km) & Norm. $r_a$ (km) & Initial $r_a$ (km) &
Nominal $h_p$ (km) \\
\midrule
Walkout & 3900 & 5000 & 4906 & 86.0 \\
Endgame & 4906 & 10100 & 6000 & 86.0 \\
Main & 3900 & 10100 & 10038 & 92.0 \\
Campaign & 3900 & 30000 & 28533 & 87.5 \\
\bottomrule
\end{tabular}
\end{center}

The default configuration is created by
\path{default_aerobraking_config}.  Its default phase is \path{Main}, its
default maximum pass count is 80, and its default backend mode is
\path{:paper_surrogate}.

\subsection{Observation Space}

The paper-compatible observation type \path{PaperObservation} contains nine
features:

\begin{center}
\begin{longtable}{>{\raggedright\arraybackslash}p{0.33\linewidth}
                  >{\raggedright\arraybackslash}p{0.18\linewidth}
                  >{\raggedright\arraybackslash}p{0.30\linewidth}}
\toprule
Feature & Units & Normalization range \\
\midrule
\endhead
\path{drag_passage_time_s} & s & $400$ to $1200$ \\
\path{apoapsis_radius_m} & m & Mars radius to phase-dependent
\path{r_norm_m} \\
\path{periapsis_altitude_m} & m & $85\,000$ to $145\,000$ \\
\path{argument_of_periapsis_rad} & rad & $0$ to $\pi$ \\
\path{raan_rad} & rad & $1.57$ to $3.14$ \\
\path{inclination_rad} & rad & $1.39$ to $1.75$ \\
\path{date_ordinal} & days & $730486$ to $731216$ \\
\path{max_density_kg_m3} & kg/m$^3$ & $0$ to $1.5\times 10^{-7}$ \\
\path{max_heat_rate_w_cm2} & W/cm$^2$ & $0$ to $0.5$ \\
\bottomrule
\end{longtable}
\end{center}

Normalization is affine:
\[
\hat{x}_i = \frac{x_i - \ell_i}{u_i - \ell_i}.
\]
The code does not clamp normalized observations; values outside the configured
bounds can therefore produce normalized values outside $[0,1]$.

\subsection{Action Space}

The action table is the 13-action paper DDQN table, in meters per second:
\[
[-1.0,\ -0.5,\ -0.3,\ -0.2,\ -0.1,\ -0.05,\ 0.0,\ 0.05,\ 0.1,\ 0.2,\ 0.3,\ 0.5,\ 1.0].
\]

Actions are represented by \path{AerobrakingAction}.  Negative actions lower
periapsis and are assigned $\phi=0^\circ$; positive actions raise periapsis and
are assigned $\phi=180^\circ$.  The zero action is index 7 in Julia's 1-based
indexing.  It has zero magnitude and is treated as no maneuver.

\subsection{Episode Reset and Randomization}

\path{reset_scenario} creates an \path{AerobrakingDecisionState}.  In nominal
mode, apoapsis, periapsis altitude, inclination, RAAN, and argument of
periapsis receive small uniform jitters around the phase constants.  In
non-nominal mode, inclination and angular elements are sampled over broader
ranges and the epoch may be shifted by up to 27 days plus an hour offset.

The randomization struct also contains process-noise fields.  When enabled,
process noise multiplicatively perturbs the surrogate apoapsis drop during
each pass.

\subsection{Current Backend Dynamics}

The active backend is \path{:paper_surrogate}.  The adapter
\path{SpaceAGORACoreAdapter} currently throws an error for any other mode.
Consequently, \path{:spaceagora_core} is not a working propagation path yet.

For each action, the surrogate transition performs the following steps:

\begin{enumerate}
\item Build a lightweight simulation-configuration named tuple containing
backend mode, phase, pass index, current apoapsis and periapsis, action
delta-V, and action angle.
\item Apply the apoapsis maneuver with a two-body energy update.  The maneuver
changes apoapsis velocity by the action magnitude and derives the resulting
periapsis altitude from orbital energy.
\item Clamp the post-maneuver periapsis altitude to $50$--$180$ km.
\item Estimate periapsis velocity, exponential Mars density, heat rate, and
drag-passage duration.
\item Compute an apoapsis drop from a base decay plus a heat-dependent decay
gain.
\item Advance the epoch by the larger of orbital period and drag-passage time.
\item Update RAAN and argument of periapsis with a J2 secular-rate
approximation.
\item Return the next decision state, observation, normalized observation,
reward, termination flags, pass metrics, and the simulation-configuration
record.
\end{enumerate}

The density model is
\[
\rho(h) = \rho_\mathrm{ref}
\exp\left(-\frac{h-h_\mathrm{ref}}{H}\right),
\]
with defaults stored in \path{AerobrakingScenarioConfig}.  The heat-rate model
scales a nominal heat rate by density and velocity:
\[
\dot{q}
= \dot{q}_\mathrm{nom}
\sqrt{\frac{\rho(h)}{\rho(h_\mathrm{nom})}}
\left(\frac{v_p}{v_\mathrm{nom}}\right)^3.
\]

\section{Reward and Termination}

\subsection{Termination}

\path{classify_termination} returns success, target undershoot, impact,
out-of-drag-passage, thermal-violation, termination, and truncation flags.
The default thresholds are:

\begin{itemize}
\item success: $|r_a-r_\mathrm{target}|\leq 10$ km;
\item target undershoot: $r_a-r_\mathrm{target}<-10$ km;
\item impact: $h_p < 85$ km;
\item out of drag passage: $h_p > 145$ km;
\item truncation: pass count at or above \path{max_passes};
\item thermal terminal: enabled in training mode when thermal status is not
nominal.
\end{itemize}

\subsection{Thermal Status}

The default heat-rate thresholds are:
\[
\dot{q}_\mathrm{low}=0.05,\quad
\dot{q}_\mathrm{high}=0.25,\quad
\dot{q}_\mathrm{medium}=0.30,\quad
\dot{q}_\mathrm{hard}=0.45
\quad \mathrm{W/cm^2}.
\]

Below the low threshold is considered a thermal violation only when the
spacecraft is still farther than the near-target band, which defaults to
100 km.  Above the high threshold is always classified as high, medium, or
hard depending on the threshold crossed.

\subsection{Reward}

The reward first applies terminal terms:

\begin{itemize}
\item impact: $-6$;
\item out of drag passage: $-6$;
\item target undershoot: $-4$;
\item success:
\[
10 - \left\lfloor \frac{|r_a-r_\mathrm{target}|}{1000}\right\rfloor
\frac{2}{5}.
\]
\end{itemize}

Thermal violations add a penalty:
\[
r_\mathrm{thermal} =
\begin{cases}
-6, & \mathrm{hard},\\
-5, & \mathrm{medium},\\
-4, & \mathrm{low\ or\ high}.
\end{cases}
\]

When the thermal status is nominal and the spacecraft has not undershot the
target, the running reward is
\[
r_\mathrm{run}
= (1-d^{0.4})
\frac{0.025}{\sigma\sqrt{2\pi}}
\exp\left[-\frac{1}{2}
\left(\frac{c_q-0.5}{\sigma}\right)^2\right]
-0.1 - 0.1\,c_{\Delta v},
\]
where $\sigma=0.1$,
\[
d = \frac{r_a-r_\mathrm{target}}
         {r_\mathrm{norm}-r_\mathrm{target}},
\qquad
c_q = \frac{\dot{q}-0.05}{0.25-0.05}.
\]
The action cost is normally $c_{\Delta v}=|\Delta v|$.  If the heat-rate cap
is below zero, the implementation sets $c_q=0$ and uses
$c_{\Delta v}=||\Delta v|-1.5|$, encouraging stronger lowering maneuvers in
very low heat-rate conditions.

\section{Baseline Policies}

Four baseline or non-learning policies are implemented:

\begin{itemize}
\item \path{NoManeuverPolicy}: always returns the zero-action index.
\item \path{RandomActionPolicy}: uniformly samples one of the 13 actions.
\item \path{FixedCorridorPolicy}: uses a fixed lowering action for low heat,
a fixed raising action for high heat, and zero action inside the nominal
corridor.
\item \path{AADSHeuristicPolicy}: predicts the next heat rate for candidate
actions and chooses the action whose prediction is closest to the target heat
rate, defaulting to $0.15$ W/cm$^2$.
\end{itemize}

\section{DDQN Implementation}

\subsection{Network}

The Q-network is implemented directly as Julia arrays rather than through an
automatic differentiation library.  The default architecture is:
\[
9 \rightarrow 1024 \rightarrow 1024 \rightarrow 13,
\]
with ReLU activations after the first two affine layers.  Weights are
orthogonally initialized and biases are initialized to zero.

The forward pass computes:
\[
\begin{aligned}
z_1 &= W_1 x + b_1, & a_1 &= \max(z_1,0),\\
z_2 &= W_2 a_1 + b_2, & a_2 &= \max(z_2,0),\\
Q(x,\cdot) &= W_3 a_2 + b_3.
\end{aligned}
\]

The loss is mean squared TD error over the selected actions:
\[
L = \frac{1}{B}\sum_{i=1}^{B}
\left(Q_\theta(s_i,a_i)-y_i\right)^2.
\]
Gradients are hand-coded through the two ReLU layers.  Gradients are clipped
by global norm before the Adam update.

\subsection{Replay Buffer}

\path{ReplayBuffer} stores fixed-capacity column-major matrices for current
and next observations plus vectors for actions, rewards, termination flags,
truncation flags, and info indices.  Insertion is circular.  Sampling is
uniform with replacement from the filled portion of the buffer.

\subsection{Target Calculation}

The learner uses Double DQN target selection.  For each sampled transition,
the online network selects the next action and the target network evaluates
that action:
\[
a^\star_i = \arg\max_a Q_\theta(s'_i,a),
\]
\[
y_i = r_i + \gamma\,\mathbb{1}_{\mathrm{bootstrap}}
Q_{\bar{\theta}}(s'_i,a^\star_i).
\]

The default discount is $\gamma=0.95$.  Terminated transitions do not
bootstrap.  Truncated transitions bootstrap by default because
\path{bootstrap_truncated} defaults to true.

\subsection{Exploration and Optimization}

\path{EpsilonSchedule} is linear after a start delay:
\[
\epsilon(t) =
\begin{cases}
\epsilon_\mathrm{start}, & t \leq t_\mathrm{start},\\
\max\left(\epsilon_\mathrm{stop},
\epsilon_\mathrm{start}
+(\epsilon_\mathrm{stop}-\epsilon_\mathrm{start})\alpha\right),
& t > t_\mathrm{start},
\end{cases}
\]
where
\[
\alpha=\min\left(1,\frac{t-t_\mathrm{start}}{t_\mathrm{decay}}\right).
\]

The default DDQN hyperparameters are:

\begin{center}
\begin{tabular}{lr}
\toprule
Parameter & Default \\
\midrule
learning rate & $10^{-4}$ \\
discount & $0.95$ \\
batch size & $256$ \\
train frequency & every 5 global steps \\
train start & 10000 global steps \\
target update & every 10000 train steps \\
gradient clip norm & $10$ \\
Adam $\epsilon$ & $10^{-6}$ \\
hidden dimension & $1024$ \\
\bottomrule
\end{tabular}
\end{center}

\section{Training Runtime}

\subsection{Configuration}

Configuration is loaded from TOML and resolved into:
\path{AerobrakingScenarioConfig}, \path{DDQNConfig}, \path{EpsilonSchedule},
\path{TrainingConfig}, and \path{ReportConfig}.  The default config path is
\path{configs/aerobraking/paper_replication.toml}.

Two included configurations are notable:

\begin{itemize}
\item \path{paper_replication.toml}: a short smoke-oriented run using
\path{global_steps = 20}.
\item \path{ddqn_full.toml}: a longer paper-style run using
\path{global_steps = 1100000}, \path{episodes = 25000}, and
\path{n_workers = 16}.
\end{itemize}

The package treats positive \path{global_steps} as the primary training
budget.  One global step corresponds to one environment transition/action
step, not one full aerobraking campaign.  If \path{global_steps} is zero or
negative, training falls back to the episode budget.

\subsection{Training Session}

\path{build_training_session} constructs:

\begin{itemize}
\item a seeded \path{MersenneTwister};
\item a \path{SpaceAGORAAerobrakingBackend};
\item a \path{DDQNLearner};
\item a \path{RunManifest};
\item a run output directory.
\end{itemize}

The manifest records the run ID, UTC creation time, phase, seed, config path
and SHA-256, backend mode, action table, normalization names, reward
configuration, and DDQN configuration.

\subsection{Threaded Rollouts}

\path{train_parallel!} is the active trainer.  It uses the smaller of
\path{n_workers} and \path{Threads.nthreads()} as the active worker count.  At
the start of each worker batch, the master copies the
online Q-network as a policy snapshot.  Each worker runs one episode with that
snapshot and returns a summary plus a list of transitions.  The master then
ingests transitions centrally:

\begin{enumerate}
\item increment global step;
\item push the transition into replay;
\item call \path{maybe_train!};
\item write periodic checkpoints if configured;
\item print progress at the configured step or episode interval.
\end{enumerate}

This design keeps learner mutation centralized while allowing rollout
generation to run concurrently.  If the global-step budget is reached in the
middle of an episode, only the remaining budgeted transitions are ingested;
the incomplete episode summary is not added to the completed-summary list.

The final checkpoint is always written to
\path{checkpoint_final.jls}.  Intermediate checkpoints are named by global
step, such as \path{checkpoint_5000.jls}.

\section{Evaluation and Artifacts}

\path{evaluate_policy} runs complete episodes for any \path{AbstractPolicy}.
\path{evaluate_baselines} evaluates the four built-in baseline policies.  The
evaluation layer computes:

\begin{itemize}
\item per-episode metrics: pass count, reward, success/failure flags, thermal
violations, target error, mission duration, total delta-V, maneuver count,
peak heat rate, minimum periapsis altitude, and solver failures;
\item aggregate metrics: success rate, mean target error, mean delta-V, mean
pass count, mean thermal violations, and total solver failures;
\item pass logs: heat rate, action, apoapsis, periapsis, angular elements,
and reward at each pass.
\end{itemize}

\path{write_evaluation_artifacts} writes three CSV files:
\path{episode_metrics.csv}, \path{summary_metrics.csv}, and
\path{pass_logs.csv}.  The example script additionally creates PNG plots and a
standalone HTML orbit visualization from pass logs.

The generalization-suite helper currently returns two suites:
\path{nominal}, which reuses the provided config, and
\path{iid_randomized}, which switches to non-nominal randomized initial
conditions.

\section{Entrypoints}

The main scripts are:

\begin{itemize}
\item \path{scripts/train.jl}: load a config, build a session, and run
\path{train_parallel!}.
\item \path{scripts/run_baselines.jl}: evaluate built-in baselines and write
CSV artifacts.
\item \path{scripts/evaluate_policy.jl}: evaluate the AADS heuristic and
print aggregate metrics.
\item \path{scripts/run_generalization_suite.jl}: evaluate baselines over
the nominal and randomized suites.
\item \path{examples/aerobraking/example_SpaceAGORA_RL_MarsOdyssey.jl}:
produce example evaluation plots and orbit visualization artifacts.
\end{itemize}

\section{Current Verification Coverage}

The RL test suite currently checks:

\begin{itemize}
\item action-table size, zero-action index, and action direction metadata;
\item normalization lower and upper bounds;
\item thermal status classification and key reward cases;
\item replay-buffer wraparound and batch shape;
\item Double DQN target calculation;
\item deterministic backend rollouts for a fixed seed and fixed action
sequence.
\end{itemize}

These tests validate the core mechanics of the current surrogate-backed
implementation.  They do not yet constitute validation against full
SpaceAGORA propagation or historical mission telemetry.

\section{Current Limitations and Open Integration Points}

\begin{itemize}
\item \textbf{Backend integration:} only \path{:paper_surrogate} is supported.
Any other backend mode currently throws an unsupported-mode error.
\item \textbf{Simulation configuration:} the configuration builder returns a
lightweight named tuple.  It is useful for tracing and future integration, but
it is not yet a full SpaceAGORA simulation object.
\item \textbf{Parallelism:} training rollouts are threaded.  The
\path{Distributed} helper can add Julia workers, but the main trainer does
not currently dispatch rollouts through distributed processes.
\item \textbf{Learning stack:} the Q-network, gradients, and Adam optimizer
are implemented manually.  This keeps dependencies small, but it also means
there is no automatic differentiation, GPU support, or optimizer ecosystem
integration in the current code.
\item \textbf{Checkpoint format:} checkpoints use Julia serialization.  This
is convenient for local runs but is not a long-term stable interchange format
across incompatible struct changes.
\item \textbf{Validation scope:} the current tests cover unit-level and
deterministic-surrogate behavior.  They do not prove mission-grade physical
fidelity.
\end{itemize}

\section{Summary}

The current RL implementation is a self-contained Julia package for
paper-compatible Mars Odyssey aerobraking experiments.  It has a clear MDP
schema, a deterministic surrogate backend, baseline policies, a manual DDQN
learner, threaded rollout orchestration, manifests, checkpoints, and CSV
evaluation artifacts.  The main engineering gap is not the RL control flow;
it is the remaining connection from the backend adapter to full SpaceAGORA
propagation and validation.

\end{document}
